\section{Síntesis y ampliación}

La línea principal de esta clase no es que un solo modelo haya ganado, sino el diseño multimodal nativo que pasa gradualmente de "unificar todo" a "interfaz unificada, mecanismos especializados". Chameleon demuestra la viabilidad del fusionamiento temprano discreto con tokens discretos; Transfusion mantiene el Transformer compartido pero hace que el texto y las imágenes utilicen objetivos diferentes; MoT dispersa los parámetros según el modo; BAGEL y $\pi_{0.7}$ llevan la especialización hacia modelos de contenido unificado y control robótico. La segunda mitad advierte: la coexistencia de arquitecturas no implica transferencia positiva de capacidades, ni mucho menos inteligencia física.

\subsection{Conclusiones del seminario y restricciones del sistema}

Ya se discutieron por separado representación, objetivo, enrutamiento y transferencia; este apartado regresa a la diapositiva 056 para comprimir esas elecciones locales en conclusiones de nivel de curso. Al leer la figura, coloque la presión computacional y la especialización en la misma cadena causal.

\lecturefigure{slide-056.jpg}{Conclusiones del seminario: problemas abiertos, brecha del mundo físico, cómputo y especialización}{Diapositiva oficial página 56; clase 00:40:06--00:41:36}

\begin{knowledgebox}{Lectura de la figura: cuatro conclusiones forman una cadena de ingeniería}
Primero, el multimodal sigue siendo un campo de investigación activo y no estandarizado; segundo, el procesamiento de información digital es fuerte, pero la capacidad del mundo físico está lejos de completarse; tercero, agregar imagen/video/audio incrementa la presión en las infraestructuras de entrenamiento e inferencia; cuarto, a corto plazo pueden aparecer más modelos específicos por capacidad. La pregunta final no es "si unificar", sino cómo organizar módulos especializados en sistemas coherentes bajo cómputo asequible e interacción confiable.
\end{knowledgebox}

\teachervoice{00:40:41--00:41:32, el docente coloca computacionalmente pesado y especialización por capacidad en la misma conclusión: debido a la complejidad del espacio de problemas y los altos costos, aumentar modelos especializados es un camino razonable a corto plazo; cómo unificarse nuevamente en sistemas será tema de investigación futura.}

\begin{importantbox}{Diez conclusiones centrales de esta clase}
\begin{enumerate}
  \item Entrada multimodal, salida multimodal e inteligencia física son capas de capacidad distintas.
  \item El token es una unidad secuencial, no necesariamente un ID de vocabulario discreto.
  \item La autorregresión global puede coexistir con la difusión dentro del bloque.
  \item Chameleon demuestra que el fusionamiento temprano discreto es viable, pero también expone el cuello de botella de cuantización.
  \item Transfusion indica que un backbone compartido no requiere una pérdida compartida.
  \item La brecha de modo es motivación para la especialización, no prueba causal.
  \item MoT usa enrutamiento determinista conocido por modo; MoE usa compuerta aprendida.
  \item Un modelo disperso debe reportar simultáneamente parámetros totales/activos y cómputo emparejado.
  \item La transferencia de comprensión y generación debe verificarse con experimentos direccionales y línea base.
  \item El lenguaje actual es un andamiaje eficiente de razonamiento, pero la inteligencia física completa del mundo requiere estado cerrado, control y seguridad.
\end{enumerate}
\end{importantbox}

\subsection{Una checklist ejecutable de revisión de arquitectura}

\begin{knowledgebox}{De demostración conceptual a sistema desplegable}
\begin{enumerate}
  \item \textbf{Representación}: ¿cada modo es código discreto, embedding semántico continuo o latente generativo?
  \item \textbf{Ordenamiento}: ¿cómo se define el orden causal y los límites entre modos?
  \item \textbf{Objetivo}: ¿qué pérdida usar para cada salida, cómo normalizar y ponderar?
  \item \textbf{Enrutamiento}: ¿qué capas comparten y cuáles dispersan por modo o experto aprendido?
  \item \textbf{Contabilidad}: ¿se reportan simultáneamente parámetros totales/activos, tokens, FLOPs, pasos de muestreo y tiempo real?
  \item \textbf{Transferencia}: ¿la formación conjunta mejora realmente respecto a la línea base de tarea única emparejada?
  \item \textbf{Servicio}: ¿cómo presupuestar streaming, latencia, KV/cache, decodificador no textual y concurrencia?
  \item \textbf{Encarnación}: si entra al mundo físico, ¿hay estimación de estado, recuperación, seguridad y métrica de éxito cerrada?
\end{enumerate}
\end{knowledgebox}

\subsection{Límites del procesamiento del apéndice del deck}

El PDF oficial tiene tres páginas adicionales de Modelos de Interacción (páginas 53--55), pero la grabación completa salta directamente desde la página 52 sobre inteligencia multimodal más amplia a la página 56 de conclusiones; ni el ponente ni las preguntas y respuestas muestran o explican esas tres páginas. Siguiendo el principio fuente primero, se mantienen en \texttt{slides-images/} y la tabla de selección, pero marcadas como opcionales, sin reconstruirse como contenido de clase. Así se preserva la integridad del deck oficial sin convertir no dicho contenido en voz docente.

\begin{warningbox}{Conclusiones que no se pueden derivar de esta clase}
No se puede derivar que todos los modos finalmente compartirán una representación; no se puede derivar que MoT sea igualmente efectivo para la comprensión de imágenes; no se puede afirmar generación más segura o realista solo con menor pérdida de difusión; no se puede tomar BAGEL/$\pi_{0.7}$ como inteligencia física completa; tampoco se puede explicar la ventaja actual de razonamiento del lenguaje como algo que cualquier inteligencia deba traducir primero a texto.
\end{warningbox}

\subsection{Lecturas adicionales y snapshot del seminario}

Se recomienda leer los papers originales por variable de diseño: primero Chameleon v2 para entender el fusionamiento temprano discreto, luego Transfusion v1 contrastando el objetivo de difusión continuo; después Mind the Gap v2 y MoT v2 diferenciando geometría de representación y evidencia arquitectónica; luego LMFusion v4 viendo la extensión de la ruta de lenguaje congelada, y finalmente BAGEL v3 con $\pi_{0.7}$ v2 comparando generación unificada digital con guiado encarnado.

\begin{knowledgebox}{Snapshot del seminario: 2026-05-21}
Este apunte se congela en la versión pública antes de la fecha de clase:
\begin{itemize}
  \item Chameleon \texttt{2405.09818v2}, VQ-VAE-2 \texttt{1906.00446v1};
  \item Transfusion \texttt{2408.11039v1}, Mind the Gap \texttt{2203.02053v2};
  \item MoT \texttt{2411.04996v2}, LMFusion \texttt{2412.15188v4};
  \item BAGEL \texttt{2505.14683v3}, $\pi_{0.7}$ \texttt{2604.15483v2}.
\end{itemize}
Todas las versiones son anteriores al 2026-05-21; cambios de producto posteriores no retroalimentan hechos de clase.
\end{knowledgebox}

Finalmente, la inteligencia multimodal nativa no consiste en conectar más entradas a un LLM, ni perseguir una "arquitectura unificada". Es un conjunto de compromisos auditable: abstracción y detalle, compartido y especializado, comprensión y generación, contenido digital y acción física, calidad del modelo y costo del sistema. Un sistema verdaderamente maduro mantendrá la interfaz de control unificada mientras reconoce que los modos distintos requieren representaciones, objetivos, capacidades y métodos de verificación diferentes.

\end{document}
